% =============================================================================
% Gravitas User Manual
% -----------------------------------------------------------------------------
% Build with:  npm run manual        (tools/build-manual.mjs, engine: Tectonic)
%
% Three files included below are GENERATED and must not be hand-edited:
%   facts.tex           counts, from the catalog / manifest / test output
%   scenarios.tex       the scenario catalog, from js/data/scenarioInfo.js
%   investigations.tex  the lesson table, from the investigation manifest
%
% Regenerate them with `npm run docs:sync` and `npm run manual`. Every number
% in this document that could go stale comes from one of them; if you find
% yourself typing a count, stop and add a macro instead.
%
% On accessibility: this file sets a document language, embeds ToUnicode maps
% so the text extracts and searches correctly, builds a bookmark outline from
% the section structure, uses no images and therefore needs no alternative
% text, and colours links dark enough to pass contrast against white. It is
% not a tagged PDF: the LaTeX in Tectonic's package bundle predates the
% tagging support that would be needed, and that limitation is stated in the
% Accessibility section rather than hidden.
% =============================================================================

\documentclass[11pt,a4paper]{article}

\usepackage{fontspec}
\setmainfont{lmroman10-regular.otf}[
  BoldFont       = lmroman10-bold.otf,
  ItalicFont     = lmroman10-italic.otf,
  BoldItalicFont = lmroman10-bolditalic.otf]
\setsansfont{lmsans10-regular.otf}[
  BoldFont   = lmsans10-bold.otf,
  ItalicFont = lmsans10-oblique.otf]
\setmonofont{lmmono10-regular.otf}[
  BoldFont   = lmmonolt10-bold.otf,
  ItalicFont = lmmono10-italic.otf,
  Scale      = 0.94]

\usepackage[margin=2.4cm,bottom=2.8cm]{geometry}
\usepackage{microtype}
\usepackage{booktabs}
\usepackage{longtable}
\usepackage{enumitem}
\usepackage{fancyhdr}
\usepackage{xcolor}
\usepackage{titlesec}

% Generated by tools/docs-facts.mjs - do not edit.
% Regenerate with: npm run docs:sync (add --full for the test counts).

\newcommand{\GravActivities}{3}
\newcommand{\GravActivityDocuments}{8}
\newcommand{\GravAxeRuns}{60}
\newcommand{\GravAxeSurfaces}{15}
\newcommand{\GravAxeThemes}{2}
\newcommand{\GravBuildCss}{200}
\newcommand{\GravBuildDeferredChunks}{178}
\newcommand{\GravBuildDeferredJs}{4172}
\newcommand{\GravBuildInitialDownload}{818}
\newcommand{\GravBuildStartupFiles}{52}
\newcommand{\GravBuildStartupJs}{618}
\newcommand{\GravConceptDoi}{10.5281/zenodo.22800609}
\newcommand{\GravDoi}{10.5281/zenodo.22800610}
\newcommand{\GravDurationAUniverseOfStars}{70-90 min}
\newcommand{\GravDurationBinaryStarPlanets}{40-50 min}
\newcommand{\GravDurationBlackHoles}{35-45 min}
\newcommand{\GravDurationButterflyEffect}{55-70 min}
\newcommand{\GravDurationDesignTheSchedule}{35-40 min}
\newcommand{\GravDurationDetectThisPlanet}{30-35 min}
\newcommand{\GravDurationGoldilocksQuestion}{40-50 min}
\newcommand{\GravDurationGravityAssist}{15-20 min}
\newcommand{\GravDurationHohmannTransfer}{20-25 min}
\newcommand{\GravDurationKeplersLaws}{35-45 min}
\newcommand{\GravDurationLagrangePoints}{25-30 min}
\newcommand{\GravDurationListeningToSpacetime}{75-90 min}
\newcommand{\GravDurationLivesOfStars}{80-100 min}
\newcommand{\GravDurationMissingMass}{45-60 min}
\newcommand{\GravDurationOrbitalEnergy}{35-45 min}
\newcommand{\GravDurationPowerLawGravity}{45-60 min}
\newcommand{\GravDurationRadialVelocity}{45-55 min}
\newcommand{\GravDurationRetrogradeMotion}{35-45 min}
\newcommand{\GravDurationTides}{35-45 min}
\newcommand{\GravDurationTransitPhotometry}{50-70 min}
\newcommand{\GravDurationTwelveNights}{40-50 min}
\newcommand{\GravDurationWeighingStars}{35-45 min}
\newcommand{\GravDurationWhatIsAGravitationalWave}{30-40 min}
\newcommand{\GravDurationWhenOrbitsLock}{55-70 min}
\newcommand{\GravETwoeFiles}{100}
\newcommand{\GravETwoeTests}{1294}
\newcommand{\GravGradedSteps}{385}
\newcommand{\GravInstructorDocuments}{54}
\newcommand{\GravInvestigationSteps}{683}
\newcommand{\GravInvestigations}{24}
\newcommand{\GravJestSuites}{189}
\newcommand{\GravJestTests}{6178}
\newcommand{\GravLocaleNames}{English, Español}
\newcommand{\GravLocales}{2}
\newcommand{\GravObjectives}{151}
\newcommand{\GravPhysicsChecks}{286}
\newcommand{\GravReleaseChecks}{48}
\newcommand{\GravScenarios}{59}
\newcommand{\GravSmallBodyScenarios}{31}
\newcommand{\GravStabilityScenarios}{12}
\newcommand{\GravStellarTracks}{8}
\newcommand{\GravStepsAUniverseOfStars}{37}
\newcommand{\GravStepsBinaryStarPlanets}{37}
\newcommand{\GravStepsBlackHoles}{30}
\newcommand{\GravStepsButterflyEffect}{28}
\newcommand{\GravStepsDesignTheSchedule}{15}
\newcommand{\GravStepsDetectThisPlanet}{26}
\newcommand{\GravStepsGoldilocksQuestion}{37}
\newcommand{\GravStepsGravityAssist}{22}
\newcommand{\GravStepsHohmannTransfer}{22}
\newcommand{\GravStepsKeplersLaws}{23}
\newcommand{\GravStepsLagrangePoints}{20}
\newcommand{\GravStepsListeningToSpacetime}{32}
\newcommand{\GravStepsLivesOfStars}{35}
\newcommand{\GravStepsMissingMass}{33}
\newcommand{\GravStepsOrbitalEnergy}{24}
\newcommand{\GravStepsPowerLawGravity}{21}
\newcommand{\GravStepsRadialVelocity}{38}
\newcommand{\GravStepsRetrogradeMotion}{33}
\newcommand{\GravStepsTides}{30}
\newcommand{\GravStepsTransitPhotometry}{29}
\newcommand{\GravStepsTwelveNights}{13}
\newcommand{\GravStepsWeighingStars}{36}
\newcommand{\GravStepsWhatIsAGravitationalWave}{28}
\newcommand{\GravStepsWhenOrbitsLock}{34}
\newcommand{\GravUiStrings}{4605}
\newcommand{\GravVersion}{1.0.0}


\definecolor{linkink}{RGB}{16,58,120}

\usepackage[
  colorlinks=true,
  linkcolor=linkink,
  urlcolor=linkink,
  citecolor=linkink,
  bookmarks=true,
  bookmarksnumbered=true,
  bookmarksopen=true,
  bookmarksopenlevel=1,
  unicode=true,
  pdflang=en-US,
  pdftitle={Gravitas User Manual},
  pdfauthor={Carl Ziegler},
  pdfsubject={User manual for Gravitas, an interactive astrophysics sandbox and
    astronomy teaching tool},
  pdfdisplaydoctitle=true
]{hyperref}

% The facts, in the document's own metadata. tools/build-manual.mjs reads this
% line back out of the finished PDF, which is how `npm run manual:check` can
% tell that the shipped manual still describes the shipped application without
% having to parse typeset text.
\hypersetup{pdfkeywords={scenarios=\GravScenarios; investigations=\GravInvestigations;
  steps=\GravInvestigationSteps; locales=\GravLocales}}

\urlstyle{same}
\setlist{itemsep=2pt,topsep=4pt,parsep=0pt}
\titleformat{\section}{\Large\bfseries\sffamily}{\thesection}{0.6em}{}
\titleformat{\subsection}{\large\bfseries\sffamily}{\thesubsection}{0.6em}{}

\pagestyle{fancy}
\fancyhf{}
\fancyhead[L]{\small\sffamily Gravitas User Manual}
\fancyhead[R]{\small\sffamily\thepage}
\renewcommand{\headrulewidth}{0.3pt}

\newcommand{\ui}[1]{\textsf{#1}}
\newcommand{\key}[1]{\texttt{#1}}

\begin{document}

\begin{titlepage}
\centering
\vspace*{4cm}
{\Huge\bfseries\sffamily Gravitas\par}
\vspace{0.6cm}
{\LARGE\sffamily User Manual\par}
\vspace{1.4cm}
{\large An interactive astrophysics sandbox and astronomy teaching tool
that runs entirely in the browser\par}
\vspace{2cm}
{\large\url{https://gravitas-sim.online}\par}
\vspace{0.4cm}
{\large \GravScenarios{} scenarios \quad\textbullet\quad
\GravInvestigations{} guided investigations \quad\textbullet\quad
\GravInvestigationSteps{} steps\par}
\vfill
{\large Carl Ziegler\par}
\vspace{0.3cm}
{\small Released under the MIT licence. This manual is generated from the
application's own catalog; see \texttt{manual/} in the repository.\par}
\end{titlepage}

\tableofcontents
\newpage

% =============================================================================
\section{What Gravitas is}

Gravitas is a gravity simulation you can teach with. It runs in a browser, needs
no account and no installation, and is served as static files: nothing you do in
it reaches a server, and nothing about you is recorded by it.

It is two things at once. It is a \emph{sandbox} --- place stars, planets and
black holes, give them velocities, and watch what Newtonian gravity does with
them --- and it is a set of \emph{guided investigations} that use that sandbox to
take a student from a prediction, through a measurement, to a number they can
defend.

\subsection{What it models, and what it only draws}

This distinction matters more than any other in this manual, and the application
states it in public at \url{https://gravitas-sim.online/model/}.

\begin{description}[style=nextline,leftmargin=1em]
  \item[Calculated] Newtonian gravitation between every pair of bodies, in two
  dimensions, integrated with a scheme you choose. Orbits, energies, angular
  momenta, transit depths, radial-velocity curves, astrometric wobble, tidal
  differentials and rotation curves all come out of that integration.
  \item[Approximated] Mergers are perfectly inelastic. The gravitational-wave
  inspiral is a phenomenological energy loss, not a solution of general
  relativity. The dark-matter halo is an added force term with an adjustable
  flat rotation speed.
  \item[Drawn only] Accretion-disc glow, relativistic jets, gravitational-lensing
  distortion and the bloom around bright objects are pictures. They influence
  nothing.
\end{description}

Every claim the simulation makes about a measurable quantity is checked by a
validation suite of \GravPhysicsChecks{} deterministic tests, each with a stated
tolerance and a written reason for that tolerance. The results, and the list of
things that are deliberately \emph{not} validated, are published with the
project.

% =============================================================================
\section{Getting started}

Open \url{https://gravitas-sim.online}. The simulation starts immediately; a
first visit also shows a short front door offering three ways in, which you can
dismiss and never see again.

\subsection{The screen}

\begin{description}[style=nextline,leftmargin=1em]
  \item[The simulation] fills the window. Scroll to zoom, drag to pan, click a
  body to select it.
  \item[The readout, top left] carries the live numbers: elapsed simulated time,
  the zoom, the simulation speed, how many bodies of each kind are present, and
  --- if you ask for them in Settings --- the integrator in use and how far
  energy and angular momentum have drifted since the run began.
  \item[The control rail, right] is an accordion: \ui{Scenario}, \ui{State},
  \ui{Tools}, \ui{View}, \ui{Capture} and \ui{Learn}. One section is open at a
  time.
  \item[The transport bar, bottom] pauses, steps and scrubs. The simulation
  records its own history, so the scrubber rewinds through what has already
  happened and holds on the frame you land on until you return to live.
  \item[The scale bar, bottom left] is always on the canvas, not in the page. It
  is drawn onto the simulation itself so that a screenshot carries its own
  scale.
\end{description}

\subsection{Keyboard}

\begin{longtable}{@{}ll@{}}
\toprule
\textbf{Key} & \textbf{What it does} \\
\midrule
\endhead
\key{Space} & Pause and resume \\
\key{R} & Recentre the camera and return to 1$\times$ zoom \\
\key{P} & Take a screenshot \\
\key{E} & Export the recorded run as CSV \\
\key{?} & The full shortcut list \\
\key{Esc} & Leave lecture mode, or close the open dialog \\
\bottomrule
\end{longtable}

The complete list is in the \ui{Shortcuts} dialog, which is the authoritative
copy; this table is the handful you will use in a first session.

% =============================================================================
\section{The instruments}

Gravitas is a measuring instrument before it is a display. Everything in this
section is in the \ui{Tools} section of the rail.

\subsection{Ruler, protractor, stopwatch}

The \textbf{ruler} is dragged across the scene and reads a distance in AU and in
kilometres. The \textbf{protractor} has a vertex and two arms and reads the angle
between them. Both store their handles in \emph{world} coordinates, so panning or
zooming does not change what they say --- a ruler pinned to two bodies keeps
measuring those two bodies.

The \textbf{stopwatch} runs on simulated time, so pausing pauses it and changing
the simulation speed does not change what it measures. With a body selected it
can be latched to that body's periapsis passages: each lap is then one orbit,
timed from closest approach rather than by your reaction time, and the panel
reports the mean of the laps.

\subsection{The object inspector}

Select a body and the inspector gives its mass, radius, speed, the orbital
elements of its path about whatever it is actually orbiting --- semi-major axis,
eccentricity, period, periapsis and apoapsis --- and its kinetic, potential and
total energy.

\subsection{Vectors and the potential well}

For the selected body, Gravitas can draw its velocity, its total acceleration,
and one arrow per gravitational source acting on it, each in a colour that cannot
be confused with the others. The arrows are the acceleration the integrator
actually used, not a second calculation that might disagree with it.

They exist for one misconception in particular: on an eccentric orbit velocity
and force are perpendicular at periapsis and apoapsis and never come within forty
degrees of parallel anywhere. An optional potential-well underlay shades the
whole scene by gravitational potential.

\subsection{Reference frames}

The \ui{Frame} selector re-expresses every position and every trail as a chosen
body would see them. This is not the same as following a body with the camera:
the recorded paths are redrawn. Put the Solar System into Earth's frame and Mars
traces the retrograde loop that the geocentric model was built to explain.

% =============================================================================
\section{The observing panels}

Three panels turn the simulation into an observation, and all three are driven by
the same viewing geometry. Drag the observer handle on the canvas, or use the
angle controls, and every open panel re-measures from the new direction.

\begin{description}[style=nextline,leftmargin=1em]
  \item[Transit light curve] Brightness against time for the starlight reaching
  the observer. A planet crossing its star carves a dip; the depth gives the
  planet's radius, and limb darkening is included, so the dip is not flat.
  \item[Radial velocity] The star's motion toward and away from the observer, as
  a spectrograph measures it. The wobble gives a planet mass --- times the sine
  of the inclination, which is where the $\sin i$ degeneracy stops being a
  footnote and becomes something a student can see by tilting the system.
  \item[Astrometry] The star's position on the sky over time, and the
  astrometric signature at an assumed distance.
  \item[Rotation curve] Orbital speed against distance from the centre, one point
  per body, drawn against the curve the visible mass alone would produce.
\end{description}

Open panels stack in the bottom-left corner rather than covering each other. The
spacetime view --- an optional 3-D rendering of the scene as a rubber sheet
dipping around each mass --- joins the same stack, and can be dragged out of it.

% =============================================================================
\section{Choosing the physics}

\subsection{Integrators}

Three schemes, switchable while the simulation runs, with live energy and
angular-momentum drift beside them:

\begin{description}[style=nextline,leftmargin=1em]
  \item[Symplectic Euler] The default, and what every scenario in the catalog is
  tuned against. First order, and its energy error is bounded rather than
  growing without limit.
  \item[Velocity Verlet] Second order, also symplectic. The best general choice
  when accuracy matters more than matching the shipped scenarios exactly.
  \item[RK4] Fourth order and not symplectic: locally the most accurate of the
  three and the only one whose energy error grows steadily over many orbits.
\end{description}

Each scheme's convergence order is \emph{measured} by the validation suite, not
asserted, and the tolerances of the integrated checks are derived from it.

\subsection{Dark matter}

An optional halo adds a force term with an adjustable flat rotation speed. It
changes the force law rather than the picture: switch it on with the rotation
curve open and a falling curve flattens. It stays on until you switch it off, so
a scenario loaded afterwards is running with modified gravity --- the rotation
curve panel says so.

\subsection{Tides}

Tidal stress is computed as the difference between the pull on a body's near side
and the pull on its centre. Where that difference exceeds the body's own gravity,
the body comes apart --- which is what happens in the Tidal Disruption Event
scenario, and what the Roche limit marks.

% =============================================================================
\section{Scenarios}

\GravScenarios{} scenarios ship with Gravitas. Open the gallery from
\ui{Scenario} in the rail; each is tagged by curriculum concept, so the gallery
can be scanned for the week you are teaching rather than read end to end.

Loading a scenario rebuilds the world from a seed, so the same scenario opens the
same way every time, and a share link carries the seed with it.

% Generated by tools/build-manual.mjs - do not edit.

\subsection{Binary Systems}
\begin{description}[style=nextline,leftmargin=1em]
  \item[Binary Pair: two stars, one balance point] Two stars of two solar masses each, four AU apart, circling their common center of mass once every four years. Neither one is stationary and neither one is orbiting the other: both go round the same point in between them. Watch the trails and the balance point gives itself away.
  \item[Binary Stars] A pair of suns in mutual orbit with 5 planets orbiting the binary system. The complex gravitational environment creates interesting orbital dynamics and potential habitable zones. Similar to real binary star systems like Alpha Centauri.
\end{description}

\subsection{Chaos \& Encounters}
\begin{description}[style=nextline,leftmargin=1em]
  \item[Three-Body Sensitivity Lab: the same start, twice] Three six-solar-mass stars at the corners of an equilateral triangle, rotating rigidly about their common centre. This is an exact solution of the three-body problem, found by Lagrange in 1772, and for three equal masses it is unstable: the triangle holds for a few turns and then comes apart. Built for one experiment - run it twice from starts that differ by fifteen hundred kilometres in a system a hundred and thirty million kilometres across, and watch how long the two runs stay together.
  \item[Gravity Slingshot] A massive black hole (60 M\ensuremath{\odot}) paired with a smaller companion (3 M\ensuremath{\odot}) create dramatic gravitational assists for nearby planets and gas giants. Watch objects gain tremendous velocity through close encounters, mimicking spacecraft gravity assists.
  \item[Rogue Encounter] A wandering black hole (30 M\ensuremath{\odot}) passes through a stable planetary system with 12 planets, 4 gas giants, and asteroids. Watch the dramatic orbital disruption, planet ejection, and tidal capture events as the rogue intruder wreaks havoc.
  \item[Star Frisbee] A dense stellar disk thrown past a rogue black hole. Will it be shredded or survive the flyby?
  \item[Kessler Cascade] Hundreds of micro-stars orbiting chaotically, colliding and ejecting like a debris cloud.
  \item[Alien Dyson Swarm Collapse] A hypothetical Dyson swarm of artificial satellites falls into a black hole after a catastrophic orbital failure.
  \item[Slingshot Gauntlet] A fast-moving star fired through a black hole obstacle course. Watch gravitational slingshots.
\end{description}

\subsection{Compact Objects}
\begin{description}[style=nextline,leftmargin=1em]
  \item[Black Hole Lab: a ten solar mass hole, and four things orbiting it] A single stellar-mass black hole with four bodies on stable circular orbits around it. Nothing is falling in. Gravity a long way from a black hole is the same gravity as anywhere else, and an object with sideways motion goes round it exactly as it would go round a star of the same mass. Select the black hole to read its Schwarzschild radius, its average density on that scale, its Hawking temperature and how long it has left.
  \item[GW150914: First Gravitational Wave Merger] Simulates the historic merger of two massive black holes (36 \& 29 M\ensuremath{\odot}) detected by LIGO in 2015. Watch as they spiral together, emit gravitational waves, and merge into a single, more massive black hole.
  \item[Binary Black Hole] Two stellar-mass black holes (15 \& 10 M\ensuremath{\odot}) locked in mutual orbit with spectacular relativistic jets. Watch as they spiral together, create gravitational waves, and eventually merge into a single, more massive black hole. The jets point in random directions for each black hole, creating a dynamic cosmic display.
  \item[Triple Black Hole] A chaotic three-body dance of massive black holes (20, 15, \& 10 M\ensuremath{\odot}) in a complex orbital arrangement. This unstable configuration will eventually eject one black hole while the remaining two merge. Demonstrates the chaotic nature of multi-body gravitational systems.
  \item[Supermassive Core] One enormous black hole (80 M\ensuremath{\odot}) dominates a dense stellar swarm with 50 planets, 5 gas giants, and 100 asteroids. The intense gravitational field creates spectacular accretion disks and tidal disruption events. Similar to the environment around real supermassive black holes in galactic centers.
  \item[Sagittarius A*] The Milky Way's central supermassive black hole (4000 M\ensuremath{\odot}, scaled down for simulation) with fast-moving S-stars, compact objects, and debris in extreme orbits. Witness the incredible gravitational forces and relativistic effects near our galaxy's supermassive black hole.
  \item[Neutron Star Merger] Two neutron stars (1.4 M\ensuremath{\odot} each) spiral toward each other in a death dance. This rare event produces gravitational waves, gamma-ray bursts, and creates heavy elements through r-process nucleosynthesis. Based on the LIGO-detected GW170817 event.
  \item[Pulsar with Planets] A rapidly spinning neutron star with 3 planets in tight orbits. The pulsar's intense magnetic field and radiation create a harsh environment. Based on the first confirmed exoplanets discovered around PSR B1257+12.
  \item[White Dwarf Binary] Two white dwarf stars in a close binary system with accretion between them. One star gradually steals material from its companion, potentially leading to a Type Ia supernova. Includes debris disk and stellar remnants.
  \item[Stellar Graveyard] A dynamic collection of stellar remnants: 3 black holes, 5 neutron stars, and 8 white dwarfs with surviving planets and extensive debris fields. Watch these stellar corpses interact in their final gravitational dance.
  \item[Compact Object Zoo] A diverse collection of compact objects: multiple black holes, neutron stars, and white dwarfs of various masses interacting in a dense environment. Perfect for studying the different types of stellar endpoints and their interactions.
  \item[Millisecond Pulsar] An extremely fast-spinning neutron star (recycled pulsar) with a white dwarf companion and planetary debris. These `recycled' pulsars are spun up by accretion and are among the most precise timekeepers in the universe.
  \item[Intermediate Mass BH] A rare intermediate-mass black hole (400 M\ensuremath{\odot}) in a globular cluster environment with dense stellar populations. These elusive objects bridge the gap between stellar-mass and supermassive black holes.
  \item[Micro BH Swarm] A dynamic swarm of small black holes (0.6-1.8 M\ensuremath{\odot}) with planets and gas giants in chaotic orbital dance. Watch as these stellar-mass black holes interact, merge, and create complex gravitational resonances.
  \item[Hungry Hungry Holes] Four black holes at the corners of a square, pulling stars from a shared central cluster.
  \item[Black Hole Billiards] A few small black holes orbiting a supermassive one, perturbing each other and creating chaotic motion.
\end{description}

\subsection{Exoplanets}
\begin{description}[style=nextline,leftmargin=1em]
  \item[TRAPPIST-1 System] A compact planetary system with seven Earth-sized worlds orbiting a cool red dwarf star just 40 light-years away. All planets are packed close to their tiny sun, with several in the habitable zone. Can you keep this delicate system stable?
  \item[Transit Lab: HD 209458 b] The first exoplanet ever caught crossing its star, found in 1999 after radial velocities said where to look. A hot Jupiter on a 3.5-day orbit, drawn here at true relative scale: the star is 1.155 solar radii, the planet 1.38 Jupiter radii, and the silhouette on screen is the same 12\% radius ratio the light curve reports. Open the Light Curve panel and watch the 1.7\% dip repeat.
  \item[Exoplanet Characterization Lab] HD 209458 again, but with the star free to move. In the Transit Lab it is pinned so the light curve stays centred; here both bodies orbit their common centre of mass, which is what the radial-velocity and astrometry instruments need in order to measure anything. The star circles a point 2.7 millionths of an AU away at 84 metres per second: far too small to see and easily large enough to detect. Open Radial Velocity and watch the wobble that found this planet.
  \item[Blended Binary: a hidden companion] The same star and planet as the Transit Lab, with a second star half a magnitude fainter sitting 300 AU away: far too close on the sky for a survey telescope to separate, and well inside one photometric aperture. Its light fills in part of the dip, so the transit measures shallower and the planet looks smaller than it is. Correcting for exactly this effect is what high-resolution imaging surveys of planet hosts are for.
  \item[Exoplanet Lab] A diverse collection of 120+ exoplanets, gas giants, and even pulsar planets around various stellar hosts. Explore the incredible diversity of planetary systems with interactive orbital mechanics and planetary interactions.
\end{description}

\subsection{Galaxies \& Clusters}
\begin{description}[style=nextline,leftmargin=1em]
  \item[Spiral Galaxy: what we expected] A galactic bulge with ninety stars orbiting it, each launched at exactly the speed the visible mass says it should have. This is the prediction, not the observation: with the mass concentrated in the middle, orbital speed falls away as the inverse square root of radius, the same way it does across the Solar System. Open the Rotation Curve panel and read the slope. Then load Milky Way Rotation and read it again.
  \item[Milky Way Rotation: what we actually see] The same disc, with every star moving at the same speed no matter how far out it is. That is what telescopes measure in real spiral galaxies, and it is far too fast for the stars you can see to hold on to: this scenario starts with a dark-matter halo switched on, because without one the disc does not survive. Open the Rotation Curve panel and switch the halo off to watch it come apart.
  \item[Coma Cluster: Zwicky, 1933] Twenty-four galaxies swarming in a bound cluster, on randomly oriented orbits, named after the members of the real Coma Cluster that Fritz Zwicky measured. He added up the light, added up the motions, and found the second answer hundreds of times larger than the first. He called the difference dunkle Materie and was ignored for forty years. Select a galaxy to read its speed, and work the same calculation he did.
  \item[Dense Star Cluster] A gravitationally bound collection of main-sequence stars, evolved giants, and stellar remnants with mutual gravitational interactions. Watch stellar encounters, binary formation, and the dynamic evolution of this stellar community over time.
  \item[Galactic Center] A supermassive black hole (4000 M\ensuremath{\odot}) surrounded by high-velocity stars, stellar remnants, and dense stellar populations. Experience the extreme gravitational environment with spectacular accretion, jets, and relativistic effects.
  \item[Galactic Collision] Two supermassive black holes (1.2M \& 1.0M M\ensuremath{\odot}) with hundreds of stars representing galactic cores in collision. Witness the formation of tidal streams, stellar disruption, and the eventual merger of supermassive black holes.
  \item[Quasar Cannon] A supermassive black hole is actively feeding on a dense star cluster. Watch a beam of light form as stars spiral inward.
  \item[The Pinwheel Galaxy Core] Two intermediate black holes in the center of a stellar disk. The disk forms a rotating pinwheel pattern as stars are slung around.
  \item[Tidal Arm Tango] Two black holes dance past each other, flinging stars into massive tidal arms like colliding galaxies.
\end{description}

\subsection{Habitability}
\begin{description}[style=nextline,leftmargin=1em]
  \item[Habitable Zone Lab: the inner Solar System, with the zone drawn] The Sun with Venus, Earth, Mars and Ceres on their real orbits, and the circumstellar habitable zone drawn around the star from a published prescription. Venus sits inside the inner edge and Mars outside the outer one on the conservative definition, and only one of the four has liquid water on its surface today. Switch the Habitable Zone Model setting to see the optimistic band, which reaches out past Mars.
\end{description}

\subsection{Orbits \& Kepler}
\begin{description}[style=nextline,leftmargin=1em]
  \item[Kepler's 2nd Law - Equal Areas] A planet in a nearly circular orbit and an eccentric orbiter around a central star. The area sweep visualization starts automatically for the eccentric body - watch how the wedges change shape but maintain equal area, showing why objects move faster at periapsis than at apoapsis.
\end{description}

\subsection{Orbital Resonance}
\begin{description}[style=nextline,leftmargin=1em]
  \item[Galilean Resonance: the Laplace lock] Io, Europa and Ganymede on the 4:2:1 chain Laplace explained in 1805, with Callisto outside it for contrast. The periods are not exactly 4:2:1 — Europa takes 0.37\% longer than two Io years — and that near miss is the point: what holds them is the Laplace argument, which stays near 180 degrees instead of running through every value, so the three are never all in conjunction. Callisto sits within 0.03\% of 7:3 with Ganymede and is in no resonance at all. A scale model, 100 times life size.
  \item[Broken Laplace Resonance: one percent out] The same four moons with one number changed: Europa starts one percent further from Jupiter. That is a hundred times wider than the resonance can hold, and the lock fails — the Laplace argument stops swinging about 180 degrees and starts going all the way round, once every forty-six Io orbits. Run it beside Galilean Resonance and the difference between a system that is locked and one that merely has convenient periods is on screen in under a minute.
  \item[Pluto and Neptune: the 3:2 that protects] Pluto's orbit crosses Neptune's, and they have never come close. The 3:2 resonance is why: Pluto goes round twice for every three Neptune years, and the resonant argument librates about 180 degrees rather than circulating — so every conjunction happens near Pluto's aphelion, far outside Neptune's reach. A third body runs almost the same orbit from outside the resonance; watch what becomes of it. True scale; one second is about 270 years.
  \item[Jupiter Trojans: the 1:1 co-orbital resonance] Two of the five Lagrange points are places a body can sit still relative to Jupiter forever, and thousands of asteroids do. A probe placed exactly at L4 does not move in Jupiter's rotating frame; the Trojan 617 Patroclus loops around L5 once every twelve and a half Jupiter years. A third probe starts one degree from L3 — an equilibrium too, and unstable — and leaves. The fourth runs a wider ordinary orbit and is in no resonance at all.
\end{description}

\subsection{Solar System}
\begin{description}[style=nextline,leftmargin=1em]
  \item[Solar System] A simulation of our Solar System featuring real planets with correct masses, orbital distances, diameters, and colors. Includes Mercury, Venus, Earth, Mars, Jupiter, Saturn, Uranus, Neptune with their actual properties, plus real asteroids (Ceres, Vesta, Pallas) and famous comets (Halley, Hale-Bopp, Hyakutake) with authentic orbital periods and characteristics.
  \item[Retrograde Mars: the loop that needed epicycles] The Sun, Earth and Mars at their real distances and periods, and nothing else. Watched from outside, both planets go round the Sun the same way and never turn back. Switch the reference frame to Earth, in Tools, and Mars stops circling and starts drawing a loop that doubles back on itself. Nothing about the physics changed; only the frame did. That loop is the observation Ptolemy built epicycles to reproduce and Copernicus explained away, and here you can turn it on and off with one control.
  \item[Earth-Moon System] A detailed simulation of the Earth-Moon system with accurate masses, orbital mechanics, and realistic appearances. Features Earth with its blue oceans and green continents, and the Moon with its characteristic gray surface and craters. Perfect for studying orbital dynamics and tidal effects.
  \item[Interstellar Visitor: 1I/'Oumuamua] The first object ever seen passing through the Solar System from somewhere else, on its measured orbit: perihelion inside Mercury, eccentricity 1.20, and 87.7 km/s at closest approach. Earth is shown for scale. Select the visitor and look at the sign of its total energy: it is positive, which is the whole story. It is not bound to the Sun, it was never going to stay, and it will not be back.
  \item[Kuiper Belt] An accurate simulation of our Solar System's Kuiper Belt featuring real dwarf planets (Pluto, Eris, Haumea, Makemake), large KBOs (Quaoar, Sedna, Orcus, Varuna), and smaller objects (Ixion, Huya, 2002 AW197) with realistic masses and orbital properties.
\end{description}

\subsection{Stellar Evolution}
\begin{description}[style=nextline,leftmargin=1em]
  \item[Supernova Remnant] The explosive aftermath of a massive star death: a neutron star surrounded by high-velocity debris, shocked planets, and disrupted gas giants. Experience the violent and energetic environment left behind by stellar death.
  \item[Stellar Nursery] A dense cluster of young stars around a proto-black hole. Watch interactions and ejections as the cluster evolves.
\end{description}

\subsection{Tides \& Disruption}
\begin{description}[style=nextline,leftmargin=1em]
  \item[Tidal Disruption] Multiple objects approach a supermassive black hole (2000 M\ensuremath{\odot}) and are torn apart by extreme tidal forces. Watch as planets and gas giants are stretched, disrupted, and either ejected or accreted, creating spectacular debris streams.
\end{description}


% =============================================================================
\section{Guided investigations}

\GravInvestigations{} investigations, \GravInvestigationSteps{} steps between
them, \GravGradedSteps{} of those graded, against \GravObjectives{} stated
learning objectives. Open them from \ui{Learn} in the rail.

Every investigation follows the same discipline:

\begin{enumerate}
  \item It asks for a \textbf{prediction} before it shows anything.
  \item It hands over an \textbf{instrument} and asks for a measurement.
  \item It \textbf{plots the student's own readings} back to them.
  \item It \textbf{saves progress} in the browser, so a session can be left and
  resumed.
  \item It \textbf{exports a lab report} as a PDF, which the student submits
  through whatever system the course already uses.
\end{enumerate}

Nothing is submitted anywhere by Gravitas itself. There is no server to submit
to.

% Generated by tools/build-manual.mjs - do not edit.

\begin{longtable}{@{}p{0.30\linewidth}rrp{0.26\linewidth}@{}}
\toprule
\textbf{Investigation} & \textbf{Steps} & \textbf{Graded} & \textbf{Time} \\
\midrule
\endhead
Kepler's Laws & 23 & 13 & 35-45 min \\
Why Mars Goes Backwards & 33 & 18 & 35-45 min \\
Finding Planets by Their Shadows & 29 & 13 & 50-70 min \\
Bound, Unbound and Escape & 23 & 9 & 35-45 min \\
Weighing the Stars & 35 & 17 & 35-45 min \\
Black Holes by the Numbers & 29 & 17 & 35-45 min \\
Finding Planets by Their Tug & 37 & 20 & 45-55 min \\
The Goldilocks Question & 37 & 19 & 40-50 min \\
The Missing Mass & 33 & 17 & 45-60 min \\
Tides & 30 & 16 & 35-45 min \\
The Butterfly Effect in Space & 28 & 13 & 55-70 min \\
When Orbits Lock & 33 & 16 & 55-70 min \\
Can You Detect This Planet? & 26 & 13 & 30-35 min \\
Design the Schedule & 14 & 7 & 35-40 min \\
Planets in Binary Stars & 37 & 21 & 40-50 min \\
Where Does a Gravity Assist Get Its Speed? & 16 & 8 & 15-20 min \\
Getting There From Here & 20 & 12 & 20-25 min \\
Where Can It Get To? & 18 & 7 & 20-25 min \\
\bottomrule
\end{longtable}

\subsection{Kepler's Laws}
\emph{Measure the shape, pacing and timing of real orbits} \quad 23 steps, 13 of them graded, 35-45 min, Introductory astronomy.

Work through all three of Kepler’s laws by measuring orbits rather than being shown them: find the focus of an ellipse, watch equal areas sweep out in equal times, and recover the three-halves power law by plotting it yourself.

\subsection{Why Mars Goes Backwards}
\emph{Change the frame, and fourteen centuries of epicycles fall away} \quad 33 steps, 18 of them graded, 35-45 min, Introductory astronomy.

Twice every three years Mars stops in the sky, reverses, and loops back on itself. Watched from outside, nothing of the sort happens: Earth and Mars both go round the Sun the same way and never turn back. You will measure both orbits, predict what Mars does when seen from Earth, then switch the reference frame and watch the loop draw itself. Nothing in the physics changes when you do. That is the entire point, and it is what took astronomy from Ptolemy to Copernicus.

\subsection{Finding Planets by Their Shadows}
\emph{Measure a transit, weigh what it tells you, and find what is hiding} \quad 29 steps, 13 of them graded, 50-70 min, Introductory astronomy.

Work through the transit method from first principles on HD 209458 b, the first planet ever caught crossing its star: measure a depth and turn it into a radius, correct it for limb darkening, time two transits to get a period, read an atmosphere out of the color of the dip, and finish by finding the hidden companion star that makes the planet look smaller than it is.

\subsection{Bound, Unbound and Escape}
\emph{Find out what decides whether something comes back} \quad 23 steps, 9 of them graded, 35-45 min, Introductory astronomy.

Fire something off a planet and find out what decides whether it falls back, circles forever, or leaves and never returns. Work up from the experiment to the idea behind it: every object near a star carries an amount of energy, and the sign of that one number settles the question. Finish on a real interstellar visitor and decide for yourself whether it will be back.

\subsection{Weighing the Stars}
\emph{Use an orbit to measure something you cannot put on a scale} \quad 35 steps, 17 of them graded, 35-45 min, Introductory astronomy.

Kepler’s laws end with Newton’s correction, and this is what that correction is for. Watch two stars circle each other, find the balance point they are both going round, and use nothing but the size and the timing of their orbit to work out how much each one weighs. No telescope has ever put a star on a scale; this is how it is actually done.

\subsection{Black Holes by the Numbers}
\emph{Make a black hole bigger and discover some surprising rules} \quad 29 steps, 17 of them graded, 35-45 min, Introductory astronomy.

Change one thing about a black hole, its mass, and watch four completely different properties respond. Its event horizon grows in step with the mass. Its average density falls. It gets colder. It lives dramatically longer. Two of those four surprise almost everybody, and you will predict them before you measure them.

\subsection{Finding Planets by Their Tug}
\emph{Watch a star wobble, weigh its planet, and combine the clues} \quad 37 steps, 20 of them graded, 45-55 min, Introductory astronomy.

A planet you cannot see still pulls on its star, and the star moves. Measure that motion two different ways, turn it into a mass, and combine it with the radius a transit gave you to work out what kind of world it is.

\subsection{The Goldilocks Question}
\emph{Move a planet, change its star, and decide what `habitable' really means} \quad 37 steps, 19 of them graded, 40-50 min, Introductory astronomy.

Work out for yourself why a planet twice as far from its star receives a quarter as much energy, why dim stars have their habitable zones tucked in close, and why an eccentric orbit means a planet does not receive one steady amount of light all year. Then finish with the harder question the phrase ``habitable zone'' invites people to skip: what does being inside it actually tell you?

\subsection{The Missing Mass}
\emph{Weigh a system twice, and find the two answers do not agree} \quad 33 steps, 17 of them graded, 45-60 min, Introductory astronomy.

There are two ways to weigh a system in space: add up the light, or watch how things move. For the Solar System the two agree. For a galaxy they do not, and for a cluster of galaxies they are out by more than a factor of ten. Students arrange mass and watch the rotation curve it makes, turn a measured speed into an enclosed mass, then take a real galaxy’s curve and try to fit it with stars alone — and fail, in the specific way the field failed for a decade, before adding a halo and getting it right. It closes on Zwicky’s cluster and the mass budget of the universe. It is how dark matter was found, and it is a measurement rather than a theory.

\subsection{Tides}
\emph{Stretch a world, move a moon, and discover why gravity can tear objects apart} \quad 30 steps, 16 of them graded, 35-45 min, Introductory astronomy.

Tides are not caused by strong gravity. They are caused by gravity being unequal across an object, and the whole lesson is built on that one subtraction: take the pull on the centre away from the pull on the near side and the far side, and everything from the two daily high tides to a star being shredded by a black hole falls out of what is left.

\subsection{The Butterfly Effect in Space}
\emph{Run the same system twice and find out how long the answer lasts} \quad 28 steps, 13 of them graded, 55-70 min, Introductory astronomy.

Two runs of the same three stars, started from positions differing by fifteen hundred kilometres in a system a hundred and thirty million kilometres across, end up somewhere completely different. Nothing random happens in between: the simulation is deterministic, and running it twice from exactly the same numbers gives exactly the same answer both times. Along the way you will measure a case that looks like chaos and is not, put a number on how fast prediction fails, and check that the number is a property of the physics rather than of the computer.

\subsection{When Orbits Lock}
\emph{A ratio is a hint. Find out what counts as proof} \quad 33 steps, 16 of them graded, 55-70 min, Introductory astronomy.

Three of Jupiter’s moons keep time with each other, Pluto crosses Neptune’s orbit and has never come near it, and thousands of asteroids sit sixty degrees ahead of Jupiter and stay there. All three are the same phenomenon, and none of them is explained by the thing everybody quotes: the ratio of the periods. You will measure the ratios, find that the tidiest one in the system belongs to a moon in no resonance at all, and then measure the quantity that actually settles it — an angle that either swings or goes round.

\subsection{Can You Detect This Planet?}
\emph{Two methods, the same problem: the answer was decided before the data arrived} \quad 26 steps, 13 of them graded, 30-35 min, Introductory astronomy.

A planet is either there or it is not, but whether you find it depends on choices you make before you take a single measurement. Plan two radial-velocity runs of the same star with the same instrument and the same number of nights, and find that one detects a Jupiter and the other cannot tell you anything. Then do it again with transits, where the same planet is a 587-sigma certainty from space and a 4-sigma maybe from three nights on the ground — and work out how many more nights would fix that, and where the answer stops improving.

\subsection{Design the Schedule}
\emph{Eight nights on the real instrument, and the times you choose decide the answer} \quad 14 steps, 7 of them graded, 35-40 min, Introductory astronomy.

You have eight nights and one star. Plan the run yourself in the live Radial Velocity panel, commit to a prediction, then observe two schedules side by side against the same star with the same instrument and the same noise — and watch one of them recover a Jupiter while the other cannot establish that the velocity changes at all. Then break your own result: change the seed, lose a fortnight to weather, and type a list of dates by hand, until you can say what a reported period has to carry before anybody else can check it.

\subsection{Planets in Binary Stars}
\emph{What survives around two stars, and how you would know} \quad 37 steps, 21 of them graded, 40-50 min, Introductory astronomy.

Most stars come in pairs, so most planets have to make a living in a system with two suns. Some orbits work and some do not, and the line between them is sharper than you would guess. Find it twice — once for a planet around one star, once for a planet around both — and then find out how much of what you just measured was the physics and how much was the arithmetic.

\subsection{Where Does a Gravity Assist Get Its Speed?}
\emph{The same flyby, measured in two frames, with two different answers} \quad 16 steps, 8 of them graded, 15-20 min, Introductory astronomy.

Voyager 2 arrived at Jupiter travelling ten kilometres a second and left travelling twenty-six. Jupiter did not burn any fuel for it. Fly the same manoeuvre yourself, measure it in the planet’s frame and in an inertial one, and find out why the two measurements disagree — and who actually paid.

\subsection{Getting There From Here}
\emph{Two burns, a long coast, and the arithmetic that decides both} \quad 20 steps, 12 of them graded, 20-25 min, Introductory astronomy.

A spacecraft at 1 AU, a station at 2.5 AU, and no fuel to waste. Work out both burns and the coast between them with a pencil, then fly the manoeuvre and see whether the engine agrees with you. It does — to a part in a thousand — which is what makes the two surprises in it worth trusting: you speed up to go further out, and you have to speed up again on arrival or you fall straight back.

\subsection{Where Can It Get To?}
\emph{Forbidden regions, five balance points, and one conserved number} \quad 18 steps, 7 of them graded, 20-25 min, Introductory astronomy.

Two stars on a circular orbit and a speck of dust that feels them both. There is one number you can compute about the speck that tells you where it is forbidden to be — and as you make it go faster, walls open one at a time in a fixed order. Find the five places where the speck could sit still, work out which of them it can reach, and then find out why ``can reach'' is three different questions wearing the same coat.


% =============================================================================
\section{Running a controlled experiment}

The \ui{A/B Bench}, in \ui{Tools}, is for the question a sandbox cannot
otherwise answer: \emph{what does changing this one thing actually do?}

Watching a simulation twice does not answer it, because the second run never
starts quite where the first did. The bench makes it start there exactly.

\begin{enumerate}
  \item \textbf{Capture start.} The world as it stands becomes the
  experiment's origin --- its seed, its clock, its bodies and their identities,
  the integrator, the reference frame and the observer.
  \item \textbf{Name it, and choose what to measure.} Bodies as chips,
  quantities as checkboxes. A quantity needing two bodies stays greyed out
  until two are chosen.
  \item \textbf{Record Run A.}
  \item \textbf{Return to start}, exactly. The simulation holds there.
  \item \textbf{Change one thing.}
  \item \textbf{Record Run B.}
  \item \textbf{Compare.} The two runs are drawn on one simulated-time axis,
  with a table of absolute and fractional differences.
\end{enumerate}

What can be measured: position and trajectory, separation between two bodies,
speed and velocity, distance from a chosen primary, orbital period, closest
approach, total energy, angular momentum, and the numerical drift in the last
two.

\subsection{What the bench will tell you off for}

The bench names the parameter that differed between the runs, and it says so
when more than one did --- a comparison with two independent variables cannot
say which one caused the difference. You can confirm that you meant it, because
a deliberate multivariable comparison is a legitimate thing to run.

It does not count elapsed time, the camera, the theme or a trail length as
variables. It \emph{does} count the simulation speed, because the speed
control scales the integration step and therefore changes the numerical error a
run reports, and it counts the seed and the scenario, because either one makes
it a different world.

\subsection{Honest sampling}

Two runs are never sampled at the same instants: an animation frame takes as
long as it takes, and a scenario that substeps advances its clock several times
inside one frame. So the bench records each sample with its own simulated time
and aligns the two runs on that, interpolating between samples and refusing to
extrapolate past the end of either. Where the runs do not overlap in simulated
time it says there is no comparison, rather than drawing one.

The recording indicator shows both the number of samples and the simulated
seconds they cover, because those two are not interchangeable.

\subsection{Keeping and sharing an experiment}

Experiments can be renamed, duplicated, deleted, saved in the browser and
reopened. Saved experiments are bounded --- if one is too large, or the store
is full, the bench says which and points you at the export, rather than
failing quietly.

Exporting gives two files: a \textbf{CSV} of both runs, one row per sample,
with the run named in a column and the unit in every column heading; and a
\textbf{JSON manifest} carrying the experiment's definition and its
provenance --- scenario, seed, integrator, timestep, units, a hash of the
initial state, what was measured and what changed. The manifest is the file to
paste into a lab report, and reopening it restores the experiment's setup so
somebody else can run it themselves.

\ui{Share setup} puts the same thing in a link. The link carries the starting
state and the A/B difference; it never carries the recorded runs, so what
arrives is an experiment to perform rather than somebody's results.

% =============================================================================
\section{Capturing, exporting and sharing}

\subsection{Screenshots and clips}

\ui{Capture} $\rightarrow$ \ui{Screenshot} (or \key{P}) saves a PNG with the
scenario's name, the scale bar and the elapsed simulated time drawn into the
image itself, so a figure in a lab report documents its own scale and its own
clock without a caption having to supply them.

\ui{Record Clip} records a stretch of the run to a video file: H.264 in MP4 where
the browser can encode it, WebM where it cannot. The same three facts are burned
into every frame. A recording indicator shows the elapsed time and the size, and
a take stops itself at three minutes or 80\,MB, whichever comes first.

\subsection{Data}

\ui{Export CSV} (or \key{E}) writes the recorded timeline: positions, velocities
and energies over time, plus the light curve if it is running. Every column names
its own unit --- \texttt{t\_days}, \texttt{x\_au}, \texttt{vx\_kms},
\texttt{r\_au}, \texttt{E\_kin\_J} --- and each row records what the body is
orbiting, so a period can be fitted without first having to work out what the
numbers mean. A companion notebook that reads the file ships with the project.

\subsection{Share links}

Any configuration encodes into the URL. Hand out
\texttt{gravitas-sim.online/\#\dots} as an assignment and every student opens the
identical system, down to the seed; a student can send one back as an answer.
Nothing touches a server.

\subsection{Embed mode and lecture mode}

Adding \texttt{?embed=1} to a link produces a chrome-free figure sized to its
frame, which drops into a course page or a learning-management system in an
iframe. It keeps its scale bar and its clock, and drops the panels that would not
fit.

\ui{Lecture Mode} fills the screen for projection: larger type and controls, the
Daylight theme, a spotlight pointer, and arrow keys that step through a prepared
sequence of links. \key{Esc} leaves.

% =============================================================================
\section{Languages}

The interface ships in \GravLocales{} languages --- \GravLocaleNames{} --- from a
catalog of \GravUiStrings{} strings, and all \GravInvestigations{} investigations
are translated.

A translation carries words and nothing else. No scenario name, seed, widget
identifier, numeric answer or measurement probe can be reached from a locale
file, so a mistranslation can change what a student reads but never what the
simulation measures or what counts as a correct answer.

% =============================================================================
\section{For instructors}

Instructor guides and answer keys for all \GravInvestigations{} investigations,
an adopter's guide and a curriculum map are published at
\url{https://gravitas-sim.online/instructors/}. They are generated from the
lessons themselves at build time, so an answer key cannot disagree with the
lesson it answers.

They sit behind a passphrase, and the honest description of what that means is on
the page itself: the site is static, with no server to check a credential
against, so the materials are AES-GCM encrypted at build time and decrypted in
the browser. That is real protection against casual discovery and against a
student who finds the URL. It is not protection against a determined attacker who
has the ciphertext. Instructors can request the passphrase through the contact
link on that page.

\subsection{Citing Gravitas}

If you use Gravitas in teaching or research, please cite it. The repository
carries a \texttt{CITATION.cff}, which GitHub renders as a ready-made citation,
and a Zenodo record.

% =============================================================================
\section{Limitations}

Stated plainly, because a teaching tool that hides these is a teaching tool that
will embarrass someone in front of a class:

\begin{itemize}
  \item The engine is \textbf{two-dimensional}. Inclination exists for the
  observing panels as a viewing geometry, not as a third spatial dimension the
  bodies move in.
  \item It is \textbf{Newtonian}. There is no general relativity; the
  gravitational-wave inspiral is a fitted energy loss and the innermost stable
  circular orbit is drawn, not derived.
  \item \textbf{Collisions are perfectly inelastic.} Two bodies that touch
  become one, conserving mass and momentum but not shape, spin or debris except
  where a scenario models debris deliberately.
  \item \textbf{Long runs drift.} Any integrator accumulates error. The readout
  reports that drift live rather than hiding it, and the validation suite
  measures how fast each scheme accumulates it.
  \item \textbf{Some scenarios switch off part of the model on purpose} ---
  static black holes, one-way gravity, the dark-matter halo. Those conserve less
  than the full model does, and the scenario audit names which ones and why.
\end{itemize}

% =============================================================================
\section{Accessibility}

Of the application:

\begin{itemize}
  \item The controls are real focusable elements, and the actions you need
  during a run --- pause, reset the view, screenshot, export, the shortcut list
  --- have keys of their own.
  \item The simulation canvas carries a text description, and state that only
  the picture would otherwise show is announced through a live region.
  \item Where colour carries meaning it is not carrying it alone: the run-state
  indicator shows a coloured dot \emph{and} the word, and every arrow in the
  vector overlay is named in the key beside it.
  \item Four themes, of which the default is near-black and high contrast and
  Daylight is its light counterpart. Windows high-contrast mode is honoured
  rather than overridden.
  \item Animation respects the operating system's reduced-motion setting.
\end{itemize}

Of this document: it declares its language, its text extracts and searches
correctly, its structure is a real heading hierarchy with a bookmark outline, it
contains no images and therefore nothing that needs alternative text, and its
links are coloured dark enough to hold contrast against the page.

It is \emph{not} a tagged PDF. The LaTeX distribution used to build it predates
the tagging support that PDF/UA requires. A reader who needs tagged structure is
better served by the web documentation, which is HTML.

% =============================================================================
\section{Where to find more}

\begin{description}[style=nextline,leftmargin=1em]
  \item[The application] \url{https://gravitas-sim.online}
  \item[What the model does and does not claim]
  \url{https://gravitas-sim.online/model/}
  \item[Instructor materials] \url{https://gravitas-sim.online/instructors/}
  \item[Source, issues and the technical documentation]
  \url{https://github.com/gravitas-sim/gravitas-sim.github.io}
\end{description}

\vspace{1em}
\noindent\small This manual was generated from the application's own scenario
catalog and investigation manifest. If a count in it disagrees with what you see
on screen, that is a bug: please report it.

\end{document}
